\usepackage[spanish,es-tabla]{babel}
\usepackage[utf8]{inputenc}
\usepackage[T1]{fontenc}

% --- Tipografia APA 7 --------------------------------------------------
% La clase apa7 ya cumple por defecto con margenes de 1 pulgada,
% interlineado doble y los 5 niveles de encabezado (centrado/negrita,
% izquierda/negrita, izquierda/negrita-cursiva, etc.) segun su propia
% documentacion. La 7ma edicion de APA acepta tanto fuentes serif de
% 12pt (Times New Roman) como Computer Modern 10pt (la fuente por
% defecto de LaTeX), asi que no es obligatorio cambiar la fuente. Se
% fuerza aqui Times (via mathptmx) para que el PDF final se vea
% consistente con "Times New Roman 12pt", que es el estandar mas
% reconocible y el que suele pedirse en cursos que no aceptan LaTeX.
\usepackage{mathptmx}

\usepackage{graphicx}
\usepackage{svg}

\usepackage{listings}
\usepackage{xcolor}

\lstdefinestyle{pseudocodigo}{
    basicstyle=\ttfamily\small,
    breaklines=true,
    frame=single,
    numbers=left,
    numberstyle=\tiny,
    columns=fullflexible,
    keepspaces=true
}

\usepackage{longtable}
\usepackage{booktabs}

\usepackage{hyperref}
\hypersetup{
    colorlinks=true,
    linkcolor=black,
    citecolor=black,
    urlcolor=blue
}

\usepackage[style=apa, backend=biber]{biblatex}
\DeclareLanguageMapping{spanish}{spanish-apa}
\addbibresource{referencias.bib}

% Logo institucional en el header de cada pagina (esquina superior derecha).
\IfFileExists{assets/logo-una.png}{%
    \fancyhead[R]{\includegraphics[height=1cm]{assets/logo-una.png}}%
}{}%

% Punto centralizado para overrides de encabezados (h1-h5)
% Configuracion centralizada de niveles de encabezado (h1-h5).
%
% apa7.cls YA formatea \section, \subsection, \subsubsection, \paragraph
% y \subparagraph segun las reglas de APA 7 de forma nativa y correcta:
%   h1 \section       -> centrado, negrita, mayus/minus
%   h2 \subsection    -> alineado izquierda, negrita, mayus/minus
%   h3 \subsubsection -> alineado izquierda, negrita cursiva, mayus/minus
%   h4 \paragraph     -> con sangria, negrita, termina en punto, texto sigue en la misma linea
%   h5 \subparagraph  -> con sangria, negrita cursiva, termina en punto, texto sigue en la misma linea
%
% NO tocar el formato de estos comandos con paquetes como titlesec: eso
% pisaria el comportamiento nativo de apa7.cls, que ya es compliant con el
% manual oficial (ver informe/formato.tex para el resto de paquetes usados).
%
% Si en el futuro hace falta AJUSTAR algo puntual (espaciado extra entre
% secciones, tamano de fuente especifico, etc.), centralizar esos cambios
% aqui usando \titlespacing o redefiniendo el comando puntual, en vez de
% hacerlo directo en cada archivo de seccion. Este archivo esta vacio por
% ahora a proposito: el formato nativo de apa7.cls no requiere overrides.

